% =========================
% Cas pratique 7
% =========================

\section[Cas 7]{Cas pratique 7 : Procédure de limitation}

%--------------------------
\begin{frame}{Question 1}

\begin{block}{}
L'article \textbf{L613-24} al. 1 CPI dispose que
le propriétaire du brevet peut à tout moment soit renoncer 
à la totalité du brevet ou à une ou plusieurs revendications, 
soit limiter la portée du brevet 
en modifiant une ou plusieurs revendications.\\
L'article \textbf{L613-24} al. 3 CPI dispose que
le directeur de l'Institut national de 
la propriété industrielle examine 
la conformité de la requête avec
 les dispositions réglementaires 
 mentionnées à l'alinéa précédent.
\end{block}

\begin{itemize}
  \item En conclusion, l'INPI peut statuer
  sur la partie française du brevet EP.
  L'INPI peut statuer rapidement car il s'agit
  d'un examen formel.
\end{itemize}

\end{frame}

%--------------------------
\begin{frame}{Question 2}

\begin{block}{}
Cours de cassation 30/05/2012. 
\end{block}

\begin{itemize}
  \item En conclusion, le recours de Mycos
  contre la décision du Président de l'INPI
  est irrecevable. Néanmoins, Mycos a le droit
  d'engager une action en nullité contre
  la partie française du brevet EP.
\end{itemize}

\end{frame}

%--------------------------
\begin{frame}{Question 3}

\begin{block}{}
L'article \textbf{378} CPC (sursis pour bonne
adminisation de la justice).
\end{block}

\begin{itemize}
  \item Mycos a intérêt à demander un sursis à statuer.
  En effet, l'action en contrefaçon est sans fondement
  lorsque l'action en nullité = nullité.
  \item Phytos a intérêt à demander un sursis à statuer.
  En effet, il y a monopole d'exploitation tant que
  l'action en nullité est en cours.
\end{itemize}

\end{frame}


%--------------------------
\begin{frame}{Question 4}

\begin{block}{}
\textbf{Principe de proportionnalité} (CA n° 071/2022): Enfin, le principe de proportionnalité impose au juge de s'assurer qu'il existe un juste 
équilibre entre des droits fondamentaux opposés, en l'occurrence la loyauté des preuves 
dont dépend le respect du droit au procès équitable, d'une part, et le droit de propriété des 
titulaires de droits de propriété intellectuelle qui doit leur permettre de réunir des preuves, 
dans des conditions qui ne soient pas inutilement complexes ou coûteuses, afm d'assurer 
le respect de ces droits, d'autre part. 
\end{block}

Les critères sont:
\begin{itemize}
  \item le respect du droit au procès équitable;
  \item le droit de prioriété des titulaires de droits
  de propriété intellectuelle.
\end{itemize}

\end{frame}